\section{Introduction}\label{sec:intro}

Short-term load forecasts at the level of a single household feed demand response, home energy management, local
flexibility markets and the operation of low-voltage networks that now host heat pumps, electric vehicles and rooftop
photovoltaics. Deep recurrent models, in particular the long short-term memory (LSTM) network~\cite{hochreiter1997lstm},
forecast individual households better than per-home statistical models when they are trained on many households
at once~\cite{kong2019stlf}. Pooling the readings in one place, however, is what smart-meter regulation and customer
consent increasingly discourage: half-hourly consumption reveals occupancy, routines and appliance use. Federated
learning (FL)~\cite{mcmahan2017fedavg,kairouz2021advances} offers an alternative. Each meter, home gateway or community
energy hub trains on its own data, and only model weights travel to an aggregator.

A large body of work now applies FL to residential load forecasting
\cite{fernndez2022,fekri2022,gholizadeh2022,widmer2023,qu2023,liu2023,fekri2023,das2026,wang2026,tibermacine2026}.
Two premises run through it. The first is that household data are strongly non-identically distributed (non-IID),
so that plain federated averaging (FedAvg) suffers, and that remedies such as the proximal term of
FedProx~\cite{li2020fedprox}, clustering~\cite{gholizadeh2022,dogra2023} or personalisation~\cite{qu2023,widmer2023,das2026}
are needed. The second is that FL is the communication-light, privacy-preserving option. Both premises have been tested only in
isolation. On LCL, Briggs et al.~\cite{briggs2022} found plain FedAvg worse than purely local models and recovered
the loss with clustering and local fine-tuning. Taïk and Cherkaoui~\cite{taik2020} and Savi and
Olivadese~\cite{savi2021} estimated large bandwidth savings over raw-data upload, whereas Briggs et
al.~\cite{briggs2020ccai} argued, without measurement, that weight vectors cost more bandwidth than half-hourly
readings on constrained meter links. Studies differ in dataset, cohort, horizon, model, number of rounds and evaluation
protocol, and almost none holds the households and the test set fixed while varying only how the data are split
across clients. As a result, the field cannot say how much of the accuracy lost by FL is caused by heterogeneity,
how much by splitting the data at all, and how much by the federated training protocol itself. Nor can it say
whether the bytes exchanged during training are small compared with simply uploading the raw readings.

This paper answers these questions with a controlled experiment on the Low Carbon London (LCL) smart-meter
trial~\cite{ukpn2015lcl}. We fix a stratified cohort of \NumHouseholds{} standard-tariff London households, one
chronological train/validation/test split, one LSTM of \NumParams{} parameters and one training budget. We then change
exactly one factor at a time: whether data are pooled, the number of clients the same homes are split into, how
homes are grouped into clients (from random to fully clustered by load profile), client participation and stragglers,
the aggregation rule, and the precision of the transmitted update. Every comparison uses the same test half-hours,
so differences in error are attributable to the manipulated factor alone.

\paragraph{Contributions}
\begin{enumerate}
\item \textbf{A controlled decomposition protocol for federated load forecasting.} We separate the gap between
federated and centralised training into the cost of the round protocol (one client holding all data), the cost of
fragmentation (the same homes split over $K$ clients), the cost of heterogeneity (a Dirichlet dial over real
load-profile clusters) and the cost of partial participation. Heterogeneity is reported as a \emph{measured}
index $H$ between clients, computed on training data only, rather than as a nominal partition label.
\item \textbf{Communication-priced accuracy.} Prior estimates of FL's bandwidth advantage are analytic, derived from model and data
sizes~\cite{taik2020}, or summarised as an overall ratio of exchanged data~\cite{savi2021}. We count every byte uploaded and downloaded and introduce two
budget metrics, R@$\varepsilon$ and B@$\varepsilon$: the rounds and megabytes a federation needs to come within
$\varepsilon$ of centralised validation error. We compare these with the volume of the raw training data itself.
\item \textbf{Equity of federation.} We report, household by household, who gains and who loses relative to training
alone, by consumption quintile, with paired signed-rank tests~\cite{wilcoxon1945} and Holm
correction~\cite{holm1979}. We also evaluate a zero-communication personalisation step, one epoch of local fine-tuning
of the global model.
\item \textbf{An update-dissimilarity schedule for the proximal coefficient, and why it does not help here.}
Adaptive proximal coefficients exist, driven by the training loss~\cite{li2020fedprox} or by a local-versus-global
performance index~\cite{mhou2025hapi}. HA-FedProx instead schedules $\mu$ from a parameter-free, bounded, per-round
estimate of FedProx's own $B$-dissimilarity~\cite{li2020fedprox}, computed in update space from messages the
server already receives. It is closely related to gradient diversity~\cite{yin2018diversity}. We show that on
household load data this signal tracks training progress rather than data heterogeneity, and we explain why. This
negative result applies to any scheme that adapts regularisation to observed update dissimilarity.
\item \textbf{An open, reproducible pipeline.} Preprocessing, partitioning, all arms, byte accounting, statistics,
figures and the tables of this paper are generated by one public code base from the public LCL release. A Flower
backend~\cite{beutel2020flower} reproduces the simulator exactly.
\end{enumerate}

Section~\ref{sec:related} reviews related work. Section~\ref{sec:method} defines the task, the arms, the
heterogeneity index, HA-FedProx and the communication metrics. Section~\ref{sec:setup} describes data and protocol.
Section~\ref{sec:results} reports the results, Section~\ref{sec:discussion} discusses their implications for
deployment, and Section~\ref{sec:conclusion} concludes.
